
\section{Appendix B. Implementation Details}

\subsection{B.1 The `softpick` Function}

\begin{lstlisting}
def softpick(x, dim=-1, temperature=1.0):
    orig\_dtype = x.dtype
    x = x.float()
    m = x.max(dim=dim, keepdim=True).values
    m = torch.where(m \textless{} -1e3, torch.zeros\_like(m), m)  # clamp for masked rows
    z = x - m
    exp\_z = torch.exp(z / temperature)
    exp\_neg\_m = torch.exp((-m) / temperature).clamp(max=1e4)
    p = F.relu(exp\_z - exp\_neg\_m)
    denom = p.sum(dim=dim, keepdim=True)
    zero\_mask = denom \textless{}= 1e-8
    if zero\_mask.any():
        fallback = F.softmax(x, dim=dim)
        p = torch.where(zero\_mask, fallback, p / (denom + 1e-8))
    else:
        p = p / denom
    return p.to(orig\_dtype)
\end{lstlisting}

\textbf{Note.} The clamp \texttt{max=1e4} on \texttt{exp\_neg\_m} prevents overflow. The 
\texttt{zero\_mask.any()} triggers fallback for the batch.

\subsection{B.2 Local Window Attention}

\begin{lstlisting}
def \_local\_window\_attention(self, Q, K, V, valid\_mask, is\_causal):
    B, H, L, d = Q.shape
    scale = 1.0 / math.sqrt(d)
    scores = torch.matmul(Q.float(), K.float().transpose(-2, -1)) * scale
    
    idx = torch.arange(L, device=Q.device)
    dist = (idx.unsqueeze(0) - idx.unsqueeze(1)).abs()
    mask = dist \textgreater{} self.window\_size
    if is\_causal:
        mask = mask | (idx.unsqueeze(0) \textgreater{} idx.unsqueeze(1))
    scores = scores.masked\_fill(mask.unsqueeze(0).unsqueeze(0), -10000.0)
    
    pad\_k = (~valid\_mask).unsqueeze(1).unsqueeze(2)
    scores = scores.masked\_fill(pad\_k, -10000.0)
    scores = torch.nan\_to\_num(scores, nan=-1e4, posinf=1e4, neginf=-1e4)
    
    attn = F.softmax(scores, dim=-1).to(Q.dtype)
    attn = self.dropout(attn)
    return torch.matmul(attn, V)
\end{lstlisting}

\subsection{B.3 Sparse Block Attention (Python Loop)}

\begin{lstlisting}
def \_sparse\_block\_attention(self, Q, K, V, valid\_mask, is\_causal):
    # ... pad to multiple of block\_size
    nb = Lp // self.block\_size
    
    # Compute block-level scores
    Q\_mean = Qb.mean(dim=3)
    K\_mean = Kb.mean(dim=3)
    blk\_scores = torch.matmul(Q\_mean.float(), K\_mean.float().transpose(-2, -1)) * scale
    
    if is\_causal:
        blk\_scores = blk\_scores + causal\_blk\_mask
    
    top\_k = min(self.top\_k, nb)
    \_, top\_idx = blk\_scores.topk(top\_k, dim=-1)
    if is\_causal:
        for i in range(nb):
            top\_idx[:, :, i, 0] = i
    else:
        top\_idx[:, :, :, 0] = 0
    
    out\_blocks = torch.zeros\_like(Qb)
    for i in range(nb):  # $\leftarrow$ PYTHON LOOP
        idx = top\_idx[:, :, i, :]
        idx\_exp = idx.unsqueeze(-1).unsqueeze(-1).expand(-1, -1, -1, self.block\_size, d)
        K\_sel = torch.gather(Kb, 2, idx\_exp).reshape(B, H, top\_k * self.block\_size, d)
        V\_sel = torch.gather(Vb, 2, idx\_exp).reshape(B, H, top\_k * self.block\_size, d)
        
        scores\_i = torch.matmul(Qb[:, :, i].float(), K\_sel.float().transpose(-2, -1)) * scale
        
        # Causal mask for selected tokens
        # ... 
        
        attn\_i = softpick(scores\_i, dim=-1).to(Q.dtype)
        out\_blocks[:, :, i] = torch.matmul(attn\_i, V\_sel)
    
    return out\_blocks.view(B, H, Lp, d)[:, :, :L, :]
\end{lstlisting}

\subsection{B.4 Global Spectral RFF}

\begin{lstlisting}
def \_global\_spectral\_rff\_attention(self, Q, K, V, valid\_mask, is\_causal):
    B, H, L, d = Q.shape
    pad = (~valid\_mask).unsqueeze(1).unsqueeze(-1)
    
    Qf = torch.matmul(Q.float(), self.rff\_proj.float())
    Kf = torch.matmul(K.float(), self.rff\_proj.float())
    norm = 1.0 / math.sqrt(self.rff\_features)
    phi\_Q = torch.cat([torch.cos(Qf), torch.sin(Qf)], dim=-1) * norm
    phi\_K = torch.cat([torch.cos(Kf), torch.sin(Kf)], dim=-1) * norm
    
    phi\_K = phi\_K.masked\_fill(pad, 0.0)
    V\_masked = V.float().masked\_fill(pad, 0.0)
    
    if is\_causal:
        KV = phi\_K.unsqueeze(-1) * V\_masked.unsqueeze(-2)  # [B, H, L, 2F, d]
        KV\_cumsum = torch.cumsum(KV, dim=2)
        num = torch.matmul(phi\_Q.unsqueeze(-2), KV\_cumsum).squeeze(-2)
        Z = torch.cumsum(phi\_K, dim=2)
        den = (phi\_Q * Z).sum(dim=-1, keepdim=True) + 1e-5
        return (num / den).to(Q.dtype)
    else:
        KV = torch.matmul(phi\_K.transpose(-2, -1), V\_masked)
        Z = phi\_K.sum(dim=-2, keepdim=True)
        num = torch.matmul(phi\_Q, KV)
        den = torch.matmul(phi\_Q, Z.transpose(-2, -1)) + 1e-5
        return (num / den).to(Q.dtype)
\end{lstlisting}

\textbf{The \texttt{KV} tensor is the memory bottleneck} for causal path.

\subsection{B.5 The `inject\_nasc` Exception}

\begin{lstlisting}
except Exception as e:
    logger.error("inject\_nasc failed: \%s", e)
    raise e  # $\leftarrow$ Re-raises!
\end{lstlisting}

\textbf{Contrast with other modules} that use \texttt{return None}.

\subsection{B.6 Common Pitfalls}

\subsubsection{B.6.1 Encoder Default}

Without setting \texttt{enable\_cascade\_on\_encoders=True}, encoders use dense.

\subsubsection{B.6.2 `softpick` Fallback}

The fallback is batch-level. If any row has zero denom, the whole batch 
uses softmax for that row (via \texttt{torch.where}), but the trigger is 
batch-wide.

\subsubsection{B.6.3 RFF Cumsum}

The \texttt{KV\_cumsum} tensor can be large. Consider fused kernel for long 
sequences.

\subsubsection{B.6.4 Python Loop in Sparse Block}

The \texttt{for i in range(nb)} loop is not vectorized. For small nb it's fine.

\subsubsection{B.6.5 `inject\_nasc` Raises}

Unlike \texttt{inject\_cosformer} or \texttt{inject\_fasa}, this raises. Handle with 
try/except in production.

\subsection{B.7 Diagnostics}

\subsubsection{Check Path}

\begin{lstlisting}
logger.setLevel(logging.DEBUG)
model = ...
# Should print "path=dense" or "path=cascade"
\end{lstlisting}

\subsubsection{Check Gate}

\begin{lstlisting}
summary = x.mean(dim=1)
gates = F.softmax(model.gate\_proj(summary), dim=-1)
print(f"Gate: {gates.mean(dim=0)}")
\end{lstlisting}

\subsubsection{Check RFF Memory}

\begin{lstlisting}
def measure\_rff\_memory(B, H, L, F, d):
    return B * H * L * 2 * F * d * 4 / 1e6  # MB
print(measure\_rff\_memory(8, 12, 128, 64, 64))  # ~50 MB
\end{lstlisting}
\begin{lstlisting}

---
\end{lstlisting}